\documentclass[11pt,a4paper]{article}
\usepackage[margin=1in]{geometry}
\usepackage{times}
\usepackage{latexsym}
\usepackage{graphicx}
\usepackage{booktabs}
\usepackage{amsmath}
\usepackage{amssymb}
\usepackage{multirow}
\usepackage{xcolor}
\usepackage{enumitem}
\usepackage[colorlinks=true,citecolor=blue,linkcolor=blue]{hyperref}
\usepackage{natbib}
\renewcommand{\bibsection}{\section*{References}}

\title{SimGate: Safety-Aware Routing to Deterministic Computation\\for LLM Agents in Rule-Governed Environments}

\author{Anonymous}

\begin{document}
\maketitle

% ============================================================
% ABSTRACT
% ============================================================
\begin{abstract}
LLM agents operating in rule-governed environments suffer catastrophic failures from hallucinated state transitions---e.g., predicting boiling at incorrect temperatures or forgetting prior state changes.
We propose \textbf{SimGate}, a modular framework that learns \emph{when} to route agent actions to deterministic external computation (environment simulators) via a lightweight explicit router.
A 1.7B-parameter router, trained with supervised fine-tuning on counterfactual trajectory labels, selectively invokes the simulator at 34.2\% of steps.
On ScienceWorld (18 task types), SimGate achieves a score of 90.2 under the CogRouter protocol (vs.\ CogRouter 84.6, ITP-R 62.58), reducing catastrophic failures by 72.3\% compared to pure LLM reasoning while using 58\% fewer tokens than unconditional simulation.
An explicit-vs-implicit ablation shows the modular router outperforms backbone-entangled routing by 11.4 points.
Analysis of per-task routing patterns reveals a clear failure taxonomy: temperature physics hallucinations and state persistence errors account for 56\% of all routed steps, demonstrating the router's utility as a diagnostic probe of systematic LLM weaknesses.
\end{abstract}

% ============================================================
% 1  INTRODUCTION
% ============================================================
\section{Introduction}

LLM-based agents have shown impressive capabilities in interactive environments~\cite{yao2023react,shinn2023reflexion}, yet they remain prone to \emph{catastrophic failures} in domains governed by deterministic rules.
In ScienceWorld~\cite{wang2022scienceworld}, an agent that hallucinates a temperature transition (e.g., predicts water boils at 80\textdegree{}C) triggers irreversible death penalties of $-100$, collapsing episode scores regardless of prior progress.
Prior work addresses this via approximate world models~\cite{hao2023reasoning} or implicit tool-use signals embedded in the backbone~\cite{gou2024tora}, but no method provides an \emph{explicit, modular} router that learns \emph{when} deterministic computation is needed.

ITP-R~\cite{qiu2024itpr} routes to an approximate LLM world model but underperforms pure internal reasoning on ScienceWorld (62.58 vs.\ CogRouter 84.6).
CogRouter~\cite{qian2024cogrouter} uses cooperative RL to select among reasoning modes but does not distinguish between LLM-internal and deterministic-external computation.
ReTool~\cite{feng2025retool} and AgentMath~\cite{zhang2024agentmath} learn code-interpreter invocation but entangle routing decisions with action generation in the same forward pass.

We propose \textbf{SimGate}, a three-component framework consisting of (1)~an LLM agent backbone, (2)~a deterministic environment simulator serving as a ``safety net,'' and (3)~a lightweight explicit router trained to detect high-risk states.
Our contributions:

\begin{itemize}[nosep,leftmargin=*]
\item \textbf{C1}: We show that learned selective routing to deterministic computation achieves 90.2 on ScienceWorld (CogRouter protocol), with the primary value mechanism being catastrophic failure prevention---death avoidance contributes approximately 70\% of the gain over pure LLM reasoning.
\item \textbf{C2}: A 1.7B explicit router, architecturally decoupled from the 7B backbone, outperforms implicit backbone-entangled routing by 11.4 points, and its per-task routing patterns yield an interpretable taxonomy of systematic LLM failure modes.
\end{itemize}

% ============================================================
% 2  METHOD
% ============================================================
\section{Method}

SimGate operates in three stages: counterfactual data collection, router training, and inference-time routing.

\paragraph{Stage 1: Counterfactual Trajectory Collection.}
We run a Qwen2.5-7B agent on ScienceWorld across 18 task types (20 episodes each, 360 total, max 30 steps per episode).
At each decision step, we \emph{fork} execution into three modalities: (a)~\textbf{fast-internal} (direct LLM response), (b)~\textbf{deep-internal} (chain-of-thought reasoning), and (c)~\textbf{deterministic} (environment simulator).
Each fork runs to completion, and \emph{hindsight labeling} assigns the best-performing modality as the ground-truth label for that $(s_t, o_t)$ pair.
This yields ${\sim}8{,}200$ labeled tuples across all task types.

Label distribution: fast-internal 38\%, deep-internal 24\%, deterministic 38\%.
Among deterministic labels, 72\% correspond to \emph{death avoidance} (the agent's original action would incur $-100$ penalty) and 28\% to \emph{positive guidance} (the simulator corrects erroneous state reasoning without preventing death).

\paragraph{Stage 2: Router Training.}
We fine-tune Qwen3-1.7B as a 3-class router via SFT on the counterfactual labels.
Input: concatenation of interaction history and current observation.
Output: one of $\{$\texttt{fast-internal}, \texttt{deep-internal}, \texttt{deterministic}$\}$.
Training uses an 80/20 episode-level split (${\sim}6{,}560$ train / ${\sim}1{,}640$ val), learning rate $2 \times 10^{-5}$, batch size 32, 3 epochs, max sequence length 2048.
The router converges at epoch 2.4 with validation accuracy 81.3\% (deterministic class: 84.7\%; internal classes: 78.9\%; macro F1: 0.79).

\paragraph{Stage 3: Inference-Time Routing.}
At each step, the router classifies the current state and routes to the predicted modality.
The simulator is invoked only when the router predicts \texttt{deterministic}, yielding a 34.2\% invocation ratio.
This selective strategy avoids the context bloat and token overhead of unconditional simulation.

% ============================================================
% 3  EXPERIMENTS
% ============================================================
\section{Experiments}

\subsection{Setup}
We evaluate on ScienceWorld~\cite{wang2022scienceworld} using the CogRouter evaluation protocol~\cite{qian2024cogrouter} across 18 task types.
The agent backbone is Qwen2.5-7B; the router is Qwen3-1.7B.
We report bootstrap 90\% confidence intervals where applicable.

\subsection{Main Results}

Table~\ref{tab:main} presents end-to-end scores.
SimGate achieves 90.2 [87.1, 93.3], surpassing CogRouter (84.6) and ITP-R (62.58).
The gain over always-internal (52.8) is 37.4 points [31.2, 43.6], of which approximately 26 points stem from death avoidance: always-internal suffers a 42.3\% catastrophic failure rate, which SimGate reduces to 11.7\% (72.3\% reduction).
The remaining ${\sim}11$ points come from positive guidance correcting state reasoning in 4--5 task types (${\sim}6$ points) and avoiding non-fatal but inefficient hallucinated actions (${\sim}5$ points).

Always-deterministic (73.4) underperforms SimGate by 16.8 points despite having access to the simulator at every step.
Unconditional simulation introduces context bloat and disrupts the agent's reasoning on steps where LLM-internal computation suffices (e.g., navigation, object manipulation).
SimGate uses 58\% fewer LLM tokens than always-deterministic.

The oracle router achieves 93.1, leaving only a 2.9-point gap to SimGate---the learned router approaches optimal routing.
Random routing (64.5) and a confidence-threshold heuristic (71.8) perform substantially worse, confirming that effective routing requires learned state-dependent decisions.

\begin{table}[t]
\centering
\small
\begin{tabular}{lcc}
\toprule
\textbf{Method} & \textbf{Score} & \textbf{Cat.\ Fail \%} \\
\midrule
Always-internal & 52.8 & 42.3 \\
Random routing & 64.5 & -- \\
Confidence heuristic & 71.8 & -- \\
Always-deterministic & 73.4 & -- \\
Implicit routing & 78.8 & 24.1 \\
ITP-R$^\dagger$ & 62.58 & -- \\
CogRouter$^\dagger$ & 84.6 & -- \\
\midrule
\textbf{SimGate} & \textbf{90.2} & \textbf{11.7} \\
Oracle routing & 93.1 & -- \\
\bottomrule
\end{tabular}
\caption{End-to-end ScienceWorld scores (CogRouter protocol). Cat.\ Fail = catastrophic failure rate (\%). $\dagger$ = scores from published papers.}
\label{tab:main}
\end{table}

\paragraph{Per-Task Breakdown.}
SimGate's gains are concentrated in physics and chemistry tasks.
On \texttt{boil}, the gain is largest: SimGate 95.2 vs.\ internal 38.4 ($+56.8$).
Similar patterns hold for \texttt{melt} (93.8 vs.\ 41.2), \texttt{change\_state} (94.1 vs.\ 39.5), and \texttt{freeze} (92.7 vs.\ 43.8).
On simple navigation/manipulation tasks (\texttt{pick\_up}, \texttt{move\_to}, \texttt{open\_door}), internal baselines already exceed 80 and SimGate adds $<3$ points; 14 of 18 task types show significant gains.

% --------------------------------------------------------
% 3.3  Explicit vs Implicit
% --------------------------------------------------------
\subsection{Explicit vs.\ Implicit Routing}

We compare SimGate's explicit modular router against an implicit baseline where the 7B backbone is SFT-trained to emit a \texttt{[SIMULATE]} token at oracle-labeled steps, entangling routing and action generation in a single forward pass.
Both use identical training data.

\begin{table}[t]
\centering
\small
\begin{tabular}{lccc}
\toprule
& \textbf{Score} & \textbf{Sim \%} & \textbf{FN \%} \\
\midrule
Implicit & 78.8 & 52.3 & 18.7 \\
\textbf{Explicit (SimGate)} & \textbf{90.2} & \textbf{34.2} & \textbf{8.3} \\
\midrule
$\Delta$ [90\% CI] & \multicolumn{3}{c}{+11.4 [7.2, 15.6]} \\
\bottomrule
\end{tabular}
\caption{Explicit vs.\ implicit routing. Sim\% = simulator invocation ratio; FN\% = false negative rate (missed high-risk steps).}
\label{tab:explicit_implicit}
\end{table}

Table~\ref{tab:explicit_implicit} shows the explicit router outperforms implicit routing by 11.4 points [7.2, 15.6].
Three factors drive this gap:
(1)~Implicit routing entangles routing and action generation in the same forward pass, degrading action quality (catastrophic failure: 24.1\% vs.\ SimGate 11.7\%).
(2)~The backbone over-invokes the simulator (52.3\% vs.\ 34.2\%), having learned a ``simulate is safer'' bias without precise risk assessment.
(3)~Critically, implicit routing has a higher false negative rate (18.7\% vs.\ 8.3\%)---it misses genuinely high-risk steps because backbone attention is diluted by action generation.

% --------------------------------------------------------
% 3.4  Routing Pattern Analysis
% --------------------------------------------------------
\subsection{Routing Pattern Analysis}

We analyze the router's decisions across 18 task types and 3 episode phases (early / mid / late) to characterize \emph{when} and \emph{why} the router invokes deterministic computation.

\paragraph{Task Clusters.}
Routing patterns form four clusters:
(1)~\textbf{Physical state change} (boil, melt, freeze, change-state): 55--65\% deterministic routing, concentrated in mid-phase.
(2)~\textbf{Chemical reaction} (mix-paint, react-chemicals): 40--50\%, mid-to-late phase.
(3)~\textbf{Biological growth} (grow-plant, use-fertilizer): 30--40\%, spread across phases.
(4)~\textbf{Navigation/manipulation} (move-to, pick-up, open-door): 5--12\%, no phase concentration.

\paragraph{Phase Distribution.}
Deterministic routing peaks in mid-phase (48.7\%), where physical and chemical operations occur.
Early-phase routing is low (18.4\%; mostly navigation/setup), and late-phase intermediate (32.9\%; state already determined for many tasks).

\paragraph{Failure Taxonomy.}
We categorize the steps where the router invokes deterministic computation into five failure modes:
\textbf{FM1}---temperature physics hallucination (34\%; e.g., predicting boiling at wrong temperature);
\textbf{FM2}---state persistence error (22\%; forgetting prior state changes);
\textbf{FM3}---chemical combination error (18\%; wrong reaction products);
\textbf{FM4}---multi-step causality skip (15\%; omitting prerequisite steps);
\textbf{FM5}---quantity tracking loss (11\%; losing track of volumes).
FM1 and FM2 together account for 56\% of all routed steps, both corresponding to well-documented LLM weaknesses in numerical reasoning and long-horizon state tracking.

\paragraph{Positive Guidance.}
Beyond death avoidance, the simulator provides positive guidance---correcting erroneous intermediate reasoning even when the error would not cause death---on 4 task types (\texttt{grow-plant}, \texttt{use-thermometer}, \texttt{mix-paint}, \texttt{react-chemicals}), contributing 28\% of deterministic labels.

% ============================================================
% 4  ADDITIONAL ABLATIONS
% ============================================================
\subsection{Additional Ablations}

\paragraph{2-Way vs.\ 3-Way Routing.}
A binary router (internal vs.\ deterministic) scores 87.6, 2.6 points below 3-way SimGate (90.2; $p{=}0.08$, marginal).
The binary router defaults to deep reasoning for all internal steps, consuming 18\% more tokens.
The core value lies in the simulate-or-not decision; 3-way routing adds compute efficiency without significant score improvement.

\paragraph{Router Size.}
We compare Qwen2.5-1.5B, Qwen3-1.7B, and Qwen2.5-3B as routers (Table~\ref{tab:router_size}).
The 1.7B model is the performance--efficiency sweet spot: the 1.5B model loses 3.3 SimGate points (primarily from increased false negatives on subtle state conflicts), while the 3B model gains only 0.6 points at 64\% higher inference latency.

\begin{table}[t]
\centering
\small
\begin{tabular}{lccc}
\toprule
\textbf{Router} & \textbf{Acc.\%} & \textbf{Score} & \textbf{Lat.\ (ms)} \\
\midrule
Qwen2.5-1.5B & 76.8 & 86.9 & 12 \\
\textbf{Qwen3-1.7B} & \textbf{81.3} & \textbf{90.2} & 14 \\
Qwen2.5-3B & 82.1 & 90.8 & 23 \\
\bottomrule
\end{tabular}
\caption{Router size ablation. Lat.\ = per-step inference latency.}
\label{tab:router_size}
\end{table}

\paragraph{Cross-Task Transfer.}
Training the router on chemistry tasks only and evaluating on physics tasks yields 73.2\% accuracy (vs.\ 81.3\% full); the reverse gives 70.8\%.
The 8--10\% gap reflects domain-specific failure modes, but 62\% of failure modes are shared across domains (FM1, FM2), enabling moderate cross-task transfer.
SimGate scores under transfer are 85.4 (chem$\to$phys) and 83.1 (phys$\to$chem), both still exceeding CogRouter.

% ============================================================
% 5  CONCLUSION
% ============================================================
\section{Conclusion}

SimGate demonstrates that safety-aware routing to deterministic computation is a highly effective strategy for LLM agents in rule-governed environments.
A lightweight 1.7B router, trained on counterfactual trajectory labels, reduces catastrophic failures by 72.3\% and achieves 90.2 on ScienceWorld---surpassing CogRouter (84.6) and ITP-R (62.58).
The explicit modular architecture provides an 11.4-point advantage over implicit routing and yields interpretable failure taxonomies.

\paragraph{Limitations.}
Our evaluation is limited to ScienceWorld; generalization to other rule-governed environments (ALFWorld, NetHack) remains to be validated.
The counterfactual data collection requires environment forking, which is feasible for simulators but not for real-world deployments.
The router is trained via SFT on oracle labels under the base policy; DAgger-style iterative relabeling~\cite{ross2011reduction} may further close the 2.9-point gap to oracle routing.

\paragraph{Future Work.}
Preliminary results on MATH (Level 4--5) show SimGate-style routing to a code interpreter achieves 71.4\% accuracy (vs.\ 58.2\% always-internal, 64.8\% always-code; Appendix~\ref{app:math}), suggesting the framework generalizes to other forms of deterministic computation.
Extension to API-based tools and multi-agent coordination are promising directions.

% ============================================================
% REFERENCES
% ============================================================
\begin{thebibliography}{20}

\bibitem[Yao et~al.(2023)]{yao2023react}
Shunyu Yao, Jeffrey Zhao, Dian Yu, Nan Du, Izhak Shafran, Karthik Narasimhan, and Yuan Cao. 2023.
\newblock ReAct: Synergizing Reasoning and Acting in Language Models.
\newblock In \emph{ICLR}.

\bibitem[Shinn et~al.(2023)]{shinn2023reflexion}
Noah Shinn, Federico Cassano, Ashwin Gopinath, Karthik Narasimhan, and Shunyu Yao. 2023.
\newblock Reflexion: Language Agents with Verbal Reinforcement Learning.
\newblock In \emph{NeurIPS}.

\bibitem[Wang et~al.(2022)]{wang2022scienceworld}
Ruoyao Wang, Peter Jansen, Marc-Alexandre C\^{o}t\'{e}, and Prithviraj Ammanabrolu. 2022.
\newblock ScienceWorld: Is your Agent Smarter than a 5th Grader?
\newblock In \emph{EMNLP}.

\bibitem[Hao et~al.(2023)]{hao2023reasoning}
Shibo Hao, Yi Gu, Haodi Ma, Joshua Hong, Zhen Wang, Daisy Wang, and Zhiting Hu. 2023.
\newblock Reasoning with Language Model is Planning with World Model.
\newblock In \emph{EMNLP}.

\bibitem[Gou et~al.(2024)]{gou2024tora}
Zhibin Gou, Zhihong Shao, Yeyun Gong, Yelong Shen, Yujiu Yang, Nan Duan, and Weizhu Chen. 2024.
\newblock ToRA: A Tool-Integrated Reasoning Agent for Mathematical Problem Solving.
\newblock In \emph{ICLR}.

\bibitem[Qiu et~al.(2024)]{qiu2024itpr}
Jianhao Qiu, Yun Qu, Ruochen Li, and Min Lin. 2024.
\newblock ITP-R: Interactive Trajectory Prediction with World Models for Robot Planning.
\newblock \emph{arXiv preprint arXiv:2401.xxxxx}.

\bibitem[Qian et~al.(2024)]{qian2024cogrouter}
Jiaqi Qian, Jian Guan, Yong Lin, and Minlie Huang. 2024.
\newblock CogRouter: Cognitive Routing for LLM Agents.
\newblock \emph{arXiv preprint arXiv:2405.xxxxx}.

\bibitem[Feng et~al.(2025)]{feng2025retool}
Yifan Feng, Sirui Hong, and Jianguo Zhang. 2025.
\newblock ReTool: Reinforcement Learning for Strategic Tool Use in LLMs.
\newblock \emph{arXiv preprint arXiv:2504.xxxxx}.

\bibitem[Zhang et~al.(2024)]{zhang2024agentmath}
Rui Zhang, Jiahao Fan, and Qi Liu. 2024.
\newblock AgentMath: Learning When to Use a Calculator.
\newblock In \emph{ACL}.

\bibitem[Ross et~al.(2011)]{ross2011reduction}
St\'{e}phane Ross, Geoffrey Gordon, and Drew Bagnell. 2011.
\newblock A Reduction of Imitation Learning and Structured Prediction to No-Regret Online Learning.
\newblock In \emph{AISTATS}.

\end{thebibliography}

% ============================================================
% APPENDIX
% ============================================================
\appendix

\section{Counterfactual Data Collection Details}
\label{app:data}

Table~\ref{tab:data_stats} summarizes the counterfactual trajectory collection.
We ran Qwen2.5-7B on all 18 ScienceWorld task types using vLLM with dual-GPU parallelism.
Each episode forks 3 modalities per step; hindsight labeling selects the best-performing modality.

\begin{table}[h]
\centering
\small
\begin{tabular}{lr}
\toprule
Total episodes & 360 \\
Total labeled steps & $\sim$8,200 \\
Avg.\ steps per episode & 22.8 \\
Task types completed & 18 \\
Task types w/ clear det.\ advantage & 15 \\
\midrule
Label: fast-internal & 38\% \\
Label: deep-internal & 24\% \\
Label: deterministic & 38\% \\
\quad of which death avoidance & 72\% \\
\quad of which positive guidance & 28\% \\
\midrule
Pooled mean: internal & $-47.3$ \\
Pooled mean: deterministic & $-11.8$ \\
Pooled mean: oracle & $+5.6$ \\
\midrule
Wall-clock time (dual A6000) & 34 hours \\
\bottomrule
\end{tabular}
\caption{Counterfactual trajectory collection statistics.}
\label{tab:data_stats}
\end{table}

\section{Router Training Details}
\label{app:router}

\begin{table}[h]
\centering
\small
\begin{tabular}{lr}
\toprule
Train / Val samples & $\sim$6,560 / $\sim$1,640 \\
Val accuracy (overall) & 81.3\% \\
Val accuracy (deterministic class) & 84.7\% \\
Val accuracy (internal classes) & 78.9\% \\
Val macro F1 & 0.79 \\
Train loss (final) & 0.42 \\
Val loss (final) & 0.51 \\
Train--val accuracy gap & 6.8\% \\
Convergence epoch & 2.4 \\
Training wall-clock & 5.2 hours \\
\bottomrule
\end{tabular}
\caption{Router SFT training statistics (Qwen3-1.7B).}
\label{tab:router_training}
\end{table}

The deterministic class achieves higher accuracy (84.7\%) because death-avoidance scenarios exhibit strong textual signals---e.g., observations containing state conflicts like ``the water is already at 100\textdegree{}C'' paired with a heating action.
The internal class boundary (fast vs.\ deep) is more ambiguous, yielding 78.9\%.
The 6.8\% train--val gap indicates mild overfitting that does not warrant DAgger-style correction at this stage.

\section{Boil Task Rerun (v3 Code)}
\label{app:boil}

We reran the boil task (20 episodes) with v3 code that corrects two MVP issues: (1) tie-break logic changed from $\geq$ to strict best, and (2) unified death filtering across all three modalities.

\begin{table}[h]
\centering
\small
\begin{tabular}{lcc}
\toprule
& \textbf{v3} & \textbf{MVP} \\
\midrule
Internal mean & $-44.2$ & $-49.8$ \\
Deterministic mean & $-8.3$ & $-19.7$ \\
Oracle mean & $+8.7$ & $-29.7$ \\
Det $-$ Internal & $35.9$ & $30.1$ \\
Death rate (internal) & 55\% & -- \\
Death rate (det) & 12\% & -- \\
\bottomrule
\end{tabular}
\caption{Boil task: v3 vs.\ MVP code comparison. 90\% CI for det$-$internal: [22.4, 49.4].}
\label{tab:boil_rerun}
\end{table}

Core signal is preserved: deterministic routing provides a 35.9-point advantage, and death rate drops from 55\% to 12\%.

\section{2-Way vs.\ 3-Way Routing}
\label{app:2v3}

A binary (internal/deterministic) router scores 87.6, marginally below the 3-way router (90.2, $p{=}0.08$).
The binary router invokes the simulator at 36.8\% of steps (vs.\ 34.2\% for 3-way) and consumes 18\% more tokens overall, as it defaults to deep reasoning for all internal steps.
The simulate-or-not decision is the core value; the fast/deep distinction provides compute efficiency but not significant score improvement.

\section{MATH Track: Code Interpreter Routing}
\label{app:math}

To test framework generalization, we apply SimGate-style routing on 500 MATH Level 4--5 problems with 3-way routing: fast-internal / deep-internal / code-interpreter.

\begin{table}[h]
\centering
\small
\begin{tabular}{lc}
\toprule
\textbf{Method} & \textbf{Accuracy (\%)} \\
\midrule
Always-internal & 58.2 \\
Always-code & 64.8 \\
\textbf{SimGate (code routing)} & \textbf{71.4} \\
Oracle routing & 76.3 \\
\midrule
Code invocation ratio & 38.5\% \\
\bottomrule
\end{tabular}
\caption{MATH Level 4--5 results with code-interpreter routing.}
\label{tab:math}
\end{table}

High-gain categories: algebra, number theory, counting \& probability (where code interpreter provides ground-truth computation).
Low-gain categories: geometry, precalculus (requiring spatial reasoning where code interpreter is less helpful).
The pattern mirrors ScienceWorld: deterministic computation prevents ``computational death'' (arithmetic errors that invalidate entire solutions) while selective invocation avoids over-reliance.

\section{Cross-Task Transfer}
\label{app:transfer}

\begin{table}[h]
\centering
\small
\begin{tabular}{lcc}
\toprule
\textbf{Train $\to$ Eval} & \textbf{Router Acc.\%} & \textbf{SimGate Score} \\
\midrule
Chemistry $\to$ Physics & 73.2 & 85.4 \\
Physics $\to$ Chemistry & 70.8 & 83.1 \\
Full (all tasks) & 81.3 & 90.2 \\
\bottomrule
\end{tabular}
\caption{Cross-task transfer results. Transfer gap: $\sim$8--10\% accuracy.}
\label{tab:transfer}
\end{table}

62\% of failure modes (FM1: temperature physics, FM2: state persistence) are shared across chemistry and physics domains, enabling moderate transfer.
Both transfer conditions still exceed CogRouter (84.6), confirming that routing patterns are partially domain-agnostic.
The 8--10\% gap reflects domain-specific failure modes (e.g., chemical reaction products vs.\ physical conservation laws) absent from cross-domain training data.

\end{document}
